% !TeX TXS-program:compile = txs:///pdflatex/[--shell-escape]

\documentclass[11pt, letterpaper]{article}

\usepackage{minted}
\usepackage[utf8]{inputenc}
\usepackage[T1]{fontenc}
\usepackage{lmodern}
\usepackage{graphicx}
\usepackage{longtable}
\usepackage{wrapfig}
\usepackage{rotating}
\usepackage{amsmath}
\usepackage{textcomp}
\usepackage{amssymb}
\usepackage{hyperref}
\usepackage[spanish]{babel}
\usepackage[round]{natbib}
\usepackage{subcaption}


\title{\bfseries Tarea}
\author{Ángel García Báez}
\date{\today}
\setcounter{tocdepth}{3} 

\begin{document}
	
	% Página de presentación
	\begin{titlepage}
		\centering
		\includegraphics[width=0.2\textwidth]{logo.png}\par
		\vspace{1cm}
		{\LARGE \bfseries Universidad Veracruzana \par}
		\vspace{1cm}
		{\Large Maestría en Inteligencia Artificial\par}
		\vspace{3cm}
		{\LARGE \bfseries Visión por Computadora \par}
		\vspace{1cm}
		{\Large \bfseries Tarea 2. Propuesta de implementación del filtro bilineal y del filtro por vecinos más cercanos para el re-escalado de imágenes en Julia. \par}
		\vfill
		{\Large \textit{Ángel García Báez}\par}
		\vspace{1cm}
		{\Large Profesor: Dr. Héctor Acosta Mesa\par}
		\vfill
		{\Large \today \par}
	\end{titlepage}
	
	% Página exclusiva para la tabla de contenidos
	\newpage
	\tableofcontents
	\newpage
	
	\cite{gonzalez2018digital}
		
	\bibliographystyle{apalike}  % Estilo de cita, puedes cambiarlo si lo prefieres.
	\bibliography{Biblio}         % Aquí incluyes el archivo .bib (sin extensión).
		
		
	\section{Anexos}
	
	
	
	\subsection{Implementación de los filtros}
	
\begin{itemize}
	\begin{minted}[linenos,firstnumber=1]{julia}
#### Interpolación bilineal ####

# Definir el filtro bilineal
# Parametros de entrada: 
## Imagen de entrada
## Ancho de salida
## Alto de salida
function bilineal(Imagen,Ancho,Alto) 
	# Parametros de salida de la nueva matriz #
	naltura = Ancho
	nancho = Alto
	# Crear la nueva matriz de enteros de tamaño naltura X nancho
	nueva = zeros(Int,naltura,nancho);
	#### Dmensiones originales de la matriz
	Oaltura = size(Imagen)[1]
	Oancho = size(Imagen)[2]
	#### Definir la relación entre los tamaños
	escalaX = (Oancho-1)/(nancho-1);
	escalaY = (Oaltura-1)/(naltura-1);
	#### Aplicación del escalado ####
	for i in 1:naltura
		# Segundo loop
			for j in 1:nancho
			# Coordenadas de la imagen original
			X = escalaX*(j-1) + 1 # Adelante del borde
			Y = escalaY*(i-1) + 1 # Adelante del borde
			### Datos para hacer el interpolado ###
			## Coordenadas en X ##
			x1 = floor(Int,X)
			x2 = min(x1+1,Oancho)
			## Coordenadas en Y ##
			y1 = floor(Int,Y)
			y2 = min(y1+1,Oaltura)
			# Pesos para la interpolación #
			dx = X-x1
			dy = Y-y1
			### Extraer los valores de los pixeles vecinos ####
			Q11 = Imagen[y1,x1]
			Q12 = Imagen[y2,x1]
			Q21 = Imagen[y1,x2]
			Q22 = Imagen[y2,x2]
			### Crear el pixel con la interpolación ####
			px = (1-dx)*(1-dy)*Q11 +
			dx*(1-dy)*Q21 +
			(1-dx)*dy*Q12 +
			dx*(dy)*Q22
			# Guardarlo en la matriz de salida #
			nueva[i,j] = round.(Int,px)
			end
		end
	# Devolver la matriz resultante
	# Convertir en escala de grises #
	res = clamp.(nueva./255,0,1);
	# Convertir la matriz en escala de grises
	return(Gray.(res)) 
	end

#### Interpolación del vecino más cercano ####

# Función de interpolación por vecinos más cercanos
function vecinos(Imagen, Ancho, Alto)
	# Dimensiones de la nueva imagen
	naltura = Alto
	nancho = Ancho
	# Crear la nueva matriz de enteros de tamaño naltura × nancho
	nueva = zeros(Int, naltura, nancho)
	
	# Dimensiones originales de la imagen
	Oaltura, Oancho = size(Imagen)
	
	# Relación de escalado
	escalaX = Oancho / nancho
	escalaY = Oaltura / naltura
	
	# Aplicar el escalado con vecinos más cercanos
	for i in 1:naltura
		for j in 1:nancho
			# Coordenadas en la imagen original
			X = round(Int, escalaX * (j - 0.5))  # Píxel más cercano en X
			Y = round(Int, escalaY * (i - 0.5))  # Píxel más cercano en Y
			
			# Asegurar que las coordenadas estén dentro del rango
			X = clamp(X, 1, Oancho)
			Y = clamp(Y, 1, Oaltura)
			
			# Asignar el valor del píxel más cercano
			nueva[i, j] = Imagen[Y, X]
		end
	end
	
	# Convertir en escala de grises y normalizar
	res = clamp.(nueva ./ 255, 0, 1)
	return Gray.(res)
end


#### Cargar la imagen de la rosa ####
# Cargar las librerias para el procesado de imagenes
using Images
#using ImageView
# Imagen De referencia
ruta = "IMG/Fig2.19(a).jpg"
img = load(ruta)
img1 = Gray.(img)
# Guardar la imagenm
#save("IMG/IM1gray.jpg", img1)

### Pruebas para el filtro bilieal ###
#### Re escalar la imagen original a 256x256 pixeles ####

# Imagen De referencia
img1 = round.(Int,Gray.(img).* 255)
res1 = bilineal(img1,256,256)
save("IMG/b1.jpg", res1) #  Guardar el resultado
res1

#### Re escalar la imagen de 256x256 pixeles a 1024x1024 ####
# Imagen De referencia
ruta1 = "IMG/b1.jpg"
img = load(ruta1)
img2 = round.(Int,Gray.(img).* 255)
res2 = bilineal(img2,1024,1024)
save("IMG/b2.jpg", res2) #  Guardar el resultado
res2

### Pruebas para el filtro del vecino más cercano ###
#### Re escalar la imagen original a 256x256 pixeles ####

# Imagen De referencia
img3 = round.(Int,Gray.(img).* 255)
res3 = vecinos(img3,256,256)
save("IMG/v1.jpg", res3) #  Guardar el resultado
res3

#### Re escalar la imagen de 256x256 pixeles a 1024x1024 ####
# Imagen De referencia
ruta2 = "IMG/v1.jpg"
img = load(ruta2)
img4 = round.(Int,Gray.(img).* 255)
res4 = vecinos(img4,1024,1024)
save("IMG/v2.jpg", res4) #  Guardar el resultado
res4
	\end{minted}
\end{itemize}
	
	
	
	
	
	
	
	
\end{document}

